\section{Results and their scope}

The strongest new result is a strictly smaller certified exponential
upper bound for the same generic normalized contraction problem used by
the most recent packets. Set $\lambda=\ceil{\log_2s}$ and
$E(s)=s/\sqrt\lambda$. The exact coefficients are
\begin{align}
 C_T&=\frac\pi{\ln2}
 \left(\frac49\sqrt{\frac{2000000}{15773}}
             +\frac{160\sqrt{14}}{351}\right)=\Ctime,\label{eq:headline-time}\\
 C_S&=\frac\pi{\ln2}\sqrt{\frac{2000000}{47319}}=\Cspace.
 \label{eq:headline-space}
\end{align}
Under the interface of Section~\ref{sec:interface}, time is
$(s+\mathcal L)^{O(1)}2^{(C_T+o(1))E(s)}$ and space is
$(s+\mathcal L)^{O(1)}2^{(C_S+o(1))E(s)}$.
Eventual coefficients $30.414$ and $29.467$ are certified by exact rational
interval calculations. The square-root ceiling changes no leading constant
relative to the $s/\sqrt{\log_2s}$ scale used by some other packets.

\begin{center}
\begin{tabular}{p{7.5cm}rr}
\toprule
Generic normalized certificate & Time & Space\\\midrule
Best reviewed polynomial-space branch \cite{flags}&46.184846&polynomial\\
Best reviewed cached branch \cite{treecert}&30.891019&30.753913\\
New menus, earlier three-stage schedule&30.512927&29.487537\\
New menus, final four-stage schedule&\textbf{30.413349}&\textbf{29.466030}\\
\bottomrule
\end{tabular}
\end{center}

The new cached time coefficient is about $1.5463\%$ smaller and its space
coefficient about $4.1877\%$ smaller than the strongest displayed comparable
generic cached result. These percentages compare coefficients, not observed
run times. If $U_{\rm old}$ and $U_{\rm new}$ denote the explicit
upper-bound expressions, then
\[
 \log_2\frac{U_{\rm old}}{U_{\rm new}}
 =\left(0.477670194017653\ldots+o(1)\right)E(s).
\]
This is a meaningful asymptotic improvement of the upper-bound function;
it is not a lower bound on the time consumed by the older algorithm.

\paragraph{The new idea.}
The earlier best cached proof selected among three contraction trees and
enumerated all tuples of their operation supports. We replace that small
unstructured menu with a finite geometric menu and a pairwise correlation
certificate. The nine-group construction uses the four parallel classes of
$\mathbb F_3^2$; the thirteen-group construction uses cyclic shifts of one
explicit seed. The certificate controls every feasible allocation of mode
weights, including table construction. It pays for every actual recursive
child-factor evaluation. This is a corpus-relative contribution; averaging,
finite geometry, and fast matrix multiplication are classical tools.

\paragraph{A second, local asymptotic improvement.}
The normalized Cauchy matrix need not be materialized to compute its source
linear forms. We give guarded product-tree algorithms for matrix-vector
application and its transpose. At odd interpolation nodes, the features
have a diagonal--Toeplitz--diagonal representation after known powers of
two are cleared. All source forms then require one convolution. In the
scalar ring this admits bit cost
$\widetilde O(N(p+N))$, hence $\widetilde O(Np)$ for $p\ge N$, including
the precision guards. The reviewed streamed application uses quadratic
feature work. The prototype's one-convolution specialization is faster on
all five reported fixtures; the general product-tree prototype is slower
at those sizes, and those regressions are retained.

\paragraph{The requested larger breakthrough remains open.}
The global SAT order remains
\[
 \poly_k(n,m)2^{O_k(n\log\log n/\sqrt{\log n})}
\]
in the unrestricted fixed-width case and
$2^{O_{k,D}(n/\sqrt{\log n})}\poly(n,m)$ at fixed clause density.
No $n/\log n$ algorithm is proved here. The complete deliverable separates
the new proved finite lemmas, imported analytic and arithmetic ingredients,
tested implementation, and genuinely conditional ideas. No end-to-end SAT
solver or new Lean formalization is claimed.

The coefficients above are per normalized Boolean wire $s$. They cannot be
inserted unchanged in an exponent written in SAT variables $n$. The
existing density-sensitive and low-degree branches remain in the available
portfolio; they have different premises. Nor can a numerical coefficient
be claimed to dominate an unspecified constant in the new Anthropic
big-$O$ bound. These distinctions are necessary to make the comparison
mathematically meaningful.
